\section{Introduction}
The \systemname{} architecture---an LLM-empowered non-real-time (non-RT) RAN intelligent controller (RIC) acting as a \emph{guider}, and a reinforcement-learning (RL)-empowered near-real-time (near-RT) RIC acting as the \emph{implementer}---was proposed, motivated, and evaluated in~\cite{bao2026llmhric}. That work established the hierarchy and the guidance/implementation division of labour on a transmit-power allocation use case in an integrated access and backhaul network, using a Llama-3.1-8B rApp and a DDPG xApp in simulation. The present paper is a companion implementation artifact to~\cite{bao2026llmhric}. It does \emph{not} re-introduce, re-derive, or re-motivate the architecture, and it makes no architectural claim of its own. It documents how that framework was extended to O-RAN network slicing---per-S-NSSAI downlink physical resource block (PRB) allocation---and how the extended system was implemented, deployed, and exercised end to end on OpenAirInterface (OAI)~\cite{oai2026} and FlexRIC~\cite{flexric2026}. Readers seeking the architectural rationale should consult~\cite{bao2026llmhric}; readers seeking to re-implement or re-run the system should find every required identifier here.

\subsection{What is new relative to the magazine paper}
The framework is unchanged, but almost everything below the framework had to be built. Five differences are load-bearing for this appendix.
\begin{itemize}
  \item \emph{A different control quantity.} The near-RT action is no longer a continuous transmit power. It is a pair of complementary integer downlink PRB budgets over two S-NSSAIs on a 106-PRB carrier, emitted as E2SM-RC RRM policy ratios and enforced by the gNB scheduler as per-slot new-data budgets.
  \item \emph{A real RAN instead of a simulation.} Control acts on a running OAI standalone system---an OAI 5G core, one band-78 RFSimulator gNB, and five nrUE processes in separate network namespaces---through genuine E2 RIC CONTROL round trips terminated by a FlexRIC near-RT RIC. Every reported latency, jitter, and acknowledgement is measured on that path.
  \item \emph{OAI and FlexRIC modifications.} Stock OAI advertises no E2SM-RC control action that changes per-slice scheduler budgets, and its downlink allocator is slice-agnostic. Both were extended: a new RC Control Service Style~2 / Action~6 handler, and an S-NSSAI-aware downlink resource-block allocator installed as the default allocation hook (Section~\ref{sec:oai}).
  \item \emph{Measurement and identity-attribution machinery.} Service-model indications do not carry a complete RNTI-to-S-NSSAI relation, so per-slice state had to be reconstructed: aligned 50~ms slice summaries derived from FlexRIC MAC, RLC, PDCP, and GTP counters, an explicit RNTI-to-UE-to-S-NSSAI mapping with validity flags, and per-transition provenance and effect-window coverage accounting that decides which transitions are admissible for learning.
  \item \emph{Algorithmic changes required for convergence under the slicing reward.} The slicing objective is a constrained one---a protected-slice throughput floor gates all throughput utility---which produces a discontinuous reward that plain DDPG did not handle. The implementer became a TD3-style deterministic actor--critic with twin critics, delayed actor updates, a behaviour-cloning teacher term, a capped guidance-fusion weight with an intent-violation fallback, and an asynchronous learner whose serving Actor is promoted only through validated snapshots.
\end{itemize}

\subsection{Deployment gaps closed}
Turning an architectural study into a running radio experiment exposed constraints that do not appear in simulation. Language inference is far too slow and too coarsely clocked to sit on the scheduling path, so guidance must be published asynchronously and versioned. Gradient updates must not run inside the near-RT decision loop, so learning must be moved to a separate process and re-entered only through atomic model swaps. A control action is meaningful only if it was acknowledged by the RAN and if the measurement window that follows it actually covers the effect interval, so transition admissibility must be decided from data provenance rather than assumed. Finally, an E2 association per 100~ms action is not viable, so actuation must be persistent. The deployment therefore preserves the independence of seven clocks---RAN slot processing, 10~ms E2 service-model reporting, 50~ms monitor summarization, 100~ms near-RT control, 1000~ms E2SM-KPM audit collection, 10~s rApp language inference, and a free-running asynchronous learner---so that none of them can stall the others.

\subsection{What the artifact contains}
This appendix contributes an artifact, not an architecture: it neither re-introduces nor re-derives the \systemname{} design of~\cite{bao2026llmhric}. Each item below is a deliverable the reader can open, and each is verifiable in the released source and released results.
\begin{enumerate}
  \item \emph{RAN extension.} The released tree contains the OAI patches themselves: an S-NSSAI-aware downlink allocator and a new E2SM-RC Style~2 / Action~6 control path that decode externally supplied per-slice minimum, maximum, and dedicated PRB ratios, validate them at the E2 boundary, and install them into MAC state under the scheduler mutex, without allowing slice budgets to displace HARQ retransmissions or MAC control elements.
  \item \emph{Near-RT data plane.} The tree contains a monitoring xApp that converts raw FlexRIC service-model counters into aligned 50~ms per-slice summaries with explicit source priority and validity masks, together with an independent E2SM-KPM collector retained as a standards-aligned audit path rather than as the control state, and the schema of the tables both write.
  \item \emph{Guidance path.} The tree contains a locally hosted, quantized Gemma rApp that summarizes three ordered 10~s state windows and an operator intent into a bounded PRB target, an A1-like versioned policy service that activates each publication atomically, and the calibration script and its recorded outputs---the capacity normalizer, the protected-slice floor, and the policy-constraint bounds the guidance is projected into.
  \item \emph{Near-RT controller and asynchronous learner.} The tree contains a 100~ms controller that maps one scalar action into an integer feasible PRB interval, actuates through a persistent Unix-domain-socket RC xApp instead of a per-action E42 association, and records every decision with its raw state, bounds, normalization statistics, and serving-Actor version; and a separate TD3-style learner whose candidate Actors are published only after finiteness, saturation, action-shift, and holdout-value gates, and are swapped in at a tick boundary.
  \item \emph{Evaluation record.} The artifact records the experiment as run: the build commands, run protocol, campaign specification, and validation gates; the result generator; and, for the three executed arms---LLM-only control, unguided learning, and LLM-guided learning---the derived per-run metrics, the generated tables and figures, and a manifest pinning the campaign-specification digest and the OAI and FlexRIC revisions. The raw run databases are not part of the artifact (Section~\ref{sec:limitations}).
\end{enumerate}

\ifFullCampaign
The complete three-scenario, three-arm, five-seed campaign is reported.
\else
The evaluation released with this version is a balanced-load, single-seed pilot: \ValidRunCount{} of the \ExpectedRunCount{} planned primary runs, covering \AvailableScenarioCount{} traffic scenario(s) and \AvailableSeedCount{} seed(s). It is reported as a demonstration that the closed loop runs and is instrumented, not as evidence of statistical superiority; no confidence interval or ranking claim is made anywhere in this paper. Its scope, and the loss of the raw run databases behind it, are stated in Section~\ref{sec:limitations}.
\fi
The source, campaign specification, and released results are reachable through a single repository macro: \codeurlprinted.
